\section*{Limitations}
\label{sec:limitations}

\begin{itemize}
\item \textbf{One dataset.} All results come from Natural Questions over Wikipedia, which contains short factoid questions. Transfer ratios may differ for multi-hop, long-form or domain-specific questions.
\item \textbf{One reader.} All answers come from Qwen2.5-7B-Instruct in 4-bit AWQ format with greedy decoding. Another reader, size or quantisation may convert hits into answers at a different rate.
\item \textbf{One hardware setting.} All cost values come from one Tesla T4 and one 4-vCPU Xeon CPU. Throughput, latency and their ranking may differ on other hardware, batch sizes or serving stacks.
\item \textbf{Single cost measurements.} Each cost value is one measurement without a confidence interval. GPU throughput varied by $-$0.7\% to $+$3.1\% between pilot and main runs.
\item \textbf{BM25 is disadvantaged.} The sub-corpus contains BM25's own top-42 non-relevant passages for every query. The BM25--dense gap is therefore larger than it would be on the full corpus and must not be generalised.
\item \textbf{Absolute scores are inflated.} The 300,000-passage sub-corpus is 11\% of the full BEIR-NQ corpus. With fewer competing passages, absolute retrieval and answer scores are expected to be higher than on the full corpus; our claims concern relative comparisons.
\item \textbf{Model-card prefixes.} Each encoder uses its own query and passage prefixes or instructions. We did not test prefix-free or alternative-instruction variants.
\item \textbf{Fixed $k$.} The main comparison uses $k = 5$ passages. The top-$k$ ablation covers only BM25, MiniLM-L6-v2 and bge-large.
\item \textbf{EM is strict.} EM counts a correct paraphrase or a longer correct answer as wrong. Token F1 and containment give the same ordering, but no metric we use judges semantic equivalence.
\item \textbf{Contamination.} The encoders and the reader may have seen Natural Questions or its Wikipedia passages in training. We cannot measure this.
\item \textbf{Joined subset.} The answer-level results cover the 2,255 of 3,452 judged queries (65.3\%) that match an NQ-open question. We found no retrieval difference between joined and unjoined queries, but the unjoined queries may differ in other ways.
\item \textbf{Exact search only.} We use exact flat search without approximate nearest-neighbour search or quantisation. Their accuracy and cost trade-offs are outside our scope.
\item \textbf{Judge unvalidated.} The faithfulness judge is a single zero-shot prompt that we have not yet validated against human labels. The pilot suggested that it may be lenient. We do not interpret its scores before human agreement is known.
\item \textbf{One human annotator.} One person labels the 200-item validation sample. We cannot report inter-annotator agreement.
\item \textbf{No dimension ablation.} We did not evaluate Matryoshka dimension truncation for R6, so our dimension-cost conclusions come only from comparing different models. Those comparisons confound dimension with architecture and training.
\item \textbf{RAGAS not run.} RAGAS was not run because its local deployment requires additional infrastructure untested on our hardware; our own single-prompt judge (Appendix~\ref{app:judge}) is the primary faithfulness measure. Its validation against human labels is pending.
\item \textbf{Pre-registration deviations.} We changed the join gate (D1) and the sub-corpus membership rule (M1) before observing outcomes, and clarified the ablation-scope rule (F4.2) before it applied (\S\ref{sec:deviations}).
\item \textbf{Execution errors.} We underestimated the S1 CPU benchmark (one hour estimated, 8\,h 24\,min measured), and the first S3a attempt failed because of a notebook step order. Neither changed any data.
\item \textbf{Power.} With 2,255 queries, the minimum detectable EM difference is 3.4\,pp at the worst-case Holm level. Smaller EM differences between dense encoders may exist but remain undetected.
\end{itemize}
